\documentclass[11pt,a4paper]{article}
\usepackage[utf8]{inputenc}
\usepackage[T1]{fontenc}
\usepackage[french]{babel}
\usepackage{lmodern}
\usepackage{geometry}
\usepackage{amsmath,amssymb,amsthm,mathtools}
\usepackage[table]{xcolor}
\usepackage{hyperref}
\usepackage{enumitem}
\usepackage{microtype}
\usepackage{fancyhdr}
\usepackage{lastpage}
\usepackage{array}
\geometry{margin=2.35cm,headheight=14pt,headsep=10pt}
\hypersetup{
  colorlinks=true,
  linkcolor=blue,
  urlcolor=blue,
  citecolor=blue,
  pdfauthor={Mourad Zéraï},
  pdftitle={Dérivation de l'équation de Bellman pour la fonction de valeur d'état}
}

\setlength{\parskip}{0.55em}
\setlength{\parindent}{0pt}
\definecolor{Accent}{RGB}{22,76,138}
\definecolor{SoftRule}{RGB}{190,200,215}
\definecolor{HeaderBg}{RGB}{246,248,252}

\newtheorem{definition}{Définition}
\newtheorem{proposition}{Proposition}
\newtheorem{remark}{Remarque}

\pagestyle{fancy}
\fancyhf{}
\fancyfoot[C]{\small \thepage/\pageref{LastPage}}
\renewcommand{\headrulewidth}{0pt}
\renewcommand{\footrulewidth}{0pt}

\begin{document}

\noindent
\setlength{\arrayrulewidth}{0.5pt}
\renewcommand{\arraystretch}{1.18}
\color{SoftRule}
\begin{tabular}{|>{\raggedright\arraybackslash}p{0.31\textwidth}|>{\centering\arraybackslash}p{0.21\textwidth}|>{\raggedleft\arraybackslash}p{0.38\textwidth}|}
\hline
\rowcolor{HeaderBg}
\color{black}\textbf{Cours Deep Reinforcement Learning} & \color{black}\textbf{Classe : 4IA} & \color{black}\textbf{Année : 2026} \\
\hline
\color{black}\textbf{Enseignant :} Mourad Zéraï, PhD & \multicolumn{2}{>{\raggedleft\arraybackslash}p{0.61\textwidth}|}{\color{black}\textbf{Institution :} ESPRIT School of Engineering} \\
\hline
\end{tabular}
\color{black}

\vspace{1.1em}

{\LARGE\bfseries\color{Accent} Dérivation de l'équation de Bellman pour la fonction de valeur d'état}\\[0.35em]
{\large Sous une politique fixée $\pi$}\\[0.9em]
{\normalsize Note de cours rigoureuse, avec justification élémentaire de chaque passage.}

\vspace{0.4em}
{\color{SoftRule}\hrule}

\section*{1. Cadre probabiliste et hypothèses}

On se place dans un processus de décision markovien (MDP) à temps discret. On note :
\begin{itemize}[leftmargin=2em]
    \item $S_t$ l'état au temps $t$ ;
    \item $A_t$ l'action choisie au temps $t$ ;
    \item $R_{t+1}$ la récompense reçue entre les temps $t$ et $t+1$ ;
    \item $\pi(a\mid s)=\mathbb P_\pi(A_t=a\mid S_t=s)$ la politique fixée ;
    \item $p(s',r\mid s,a)=\mathbb P(S_{t+1}=s',R_{t+1}=r\mid S_t=s,A_t=a)$ le noyau de transition-récompense.
\end{itemize}

On suppose que le retour actualisé est bien défini :
\[
G_t = \sum_{k=0}^{\infty} \gamma^k R_{t+k+1},
\qquad 0\le \gamma <1,
\]
ou bien que l'épisode se termine presque sûrement avec somme finie. Cette hypothèse garantit que les espérances écrites ci-dessous existent.

\begin{definition}[Fonction de valeur d'état]
Pour tout état $s$, la fonction de valeur d'état sous la politique $\pi$ est définie par
\[
v_\pi(s) \coloneqq \mathbb E_\pi[G_t\mid S_t=s].
\]
\end{definition}

\section*{2. Identité de décomposition du retour}

Par définition du retour,
\[
G_t = R_{t+1} + \gamma R_{t+2} + \gamma^2 R_{t+3} + \cdots.
\]
En factorisant $\gamma$ dans tous les termes sauf le premier, on obtient l'identité fondamentale
\[
G_t = R_{t+1} + \gamma G_{t+1}.
\]

\textbf{Justification.}
En effet,
\[
G_{t+1}=R_{t+2}+\gamma R_{t+3}+\gamma^2R_{t+4}+\cdots,
\]
donc
\[
\gamma G_{t+1}=\gamma R_{t+2}+\gamma^2R_{t+3}+\gamma^3R_{t+4}+\cdots,
\]
et en ajoutant $R_{t+1}$ on retrouve exactement le développement de $G_t$.

\section*{3. Point de départ de la dérivation}

Partons de la définition de $v_\pi$ :
\[
v_\pi(s)=\mathbb E_\pi[G_t\mid S_t=s].
\]
En utilisant l'identité précédente, on remplace $G_t$ par $R_{t+1}+\gamma G_{t+1}$ :
\[
v_\pi(s)=\mathbb E_\pi[R_{t+1}+\gamma G_{t+1}\mid S_t=s].
\]

\textbf{Formule utilisée.}
On a simplement substitué une variable aléatoire par une variable presque sûrement égale.

Ensuite, par linéarité de l'espérance conditionnelle,
\[
v_\pi(s)=\mathbb E_\pi[R_{t+1}\mid S_t=s] + \gamma\,\mathbb E_\pi[G_{t+1}\mid S_t=s].
\]

\textbf{Théorème utilisé.}
Si $X$ et $Y$ sont intégrables et $c$ est une constante, alors
\[
\mathbb E[X+cY\mid \mathcal G]=\mathbb E[X\mid \mathcal G] + c\,\mathbb E[Y\mid \mathcal G].
\]
Ici $\mathcal G=\sigma(S_t)$.

\section*{4. Introduction de l'action puis du prochain état}

Le terme délicat est $\mathbb E_\pi[G_{t+1}\mid S_t=s]$ car $G_{t+1}$ dépend de l'action choisie et de la transition vers l'état suivant. On conditionne donc d'abord par l'action, puis par l'état suivant.

\subsection*{4.1 Conditionnement par l'action}

Par la loi de l'espérance totale,
\[
\mathbb E_\pi[G_{t+1}\mid S_t=s]
=\sum_a \mathbb P_\pi(A_t=a\mid S_t=s)\,\mathbb E_\pi[G_{t+1}\mid S_t=s,A_t=a].
\]
Comme $\mathbb P_\pi(A_t=a\mid S_t=s)=\pi(a\mid s)$,
\[
\mathbb E_\pi[G_{t+1}\mid S_t=s]
=\sum_a \pi(a\mid s)\,\mathbb E_\pi[G_{t+1}\mid S_t=s,A_t=a].
\]

\textbf{Théorème utilisé.}
Loi de l'espérance totale conditionnelle, sous forme discrète :
\[
\mathbb E[X\mid Y=y]=\sum_z \mathbb P(Z=z\mid Y=y)\,\mathbb E[X\mid Y=y,Z=z].
\]

De la même façon,
\[
\mathbb E_\pi[R_{t+1}\mid S_t=s]
=\sum_a \pi(a\mid s)\,\mathbb E_\pi[R_{t+1}\mid S_t=s,A_t=a].
\]

En réinjectant ces deux expressions,
\[
v_\pi(s)
=\sum_a \pi(a\mid s)\Big(
\mathbb E_\pi[R_{t+1}\mid S_t=s,A_t=a]
+\gamma\,\mathbb E_\pi[G_{t+1}\mid S_t=s,A_t=a]
\Big).
\]

\subsection*{4.2 Conditionnement par le prochain état et la récompense}

Fixons maintenant un état $s$ et une action $a$. Toujours par la loi de l'espérance totale,
\[
\mathbb E_\pi[R_{t+1}+\gamma G_{t+1}\mid S_t=s,A_t=a]
=\sum_{s',r} p(s',r\mid s,a)\,\mathbb E_\pi[R_{t+1}+\gamma G_{t+1}\mid S_t=s,A_t=a,S_{t+1}=s',R_{t+1}=r].
\]

Or, sur l'événement $\{S_{t+1}=s',R_{t+1}=r\}$, la variable $R_{t+1}$ vaut exactement $r$. Donc
\[
\mathbb E_\pi[R_{t+1}+\gamma G_{t+1}\mid S_t=s,A_t=a,S_{t+1}=s',R_{t+1}=r]
= r + \gamma\,\mathbb E_\pi[G_{t+1}\mid S_t=s,A_t=a,S_{t+1}=s',R_{t+1}=r].
\]

\section*{5. Utilisation de la propriété de Markov}

La propriété de Markov affirme que, conditionnellement à l'état courant et à l'action courante, la loi du prochain couple $(S_{t+1},R_{t+1})$ ne dépend pas du passé. Pour la suite du processus, une fois $S_{t+1}$ connu, l'avenir ne dépend plus du passé que par cet état. Ainsi,
\[
\mathbb E_\pi[G_{t+1}\mid S_t=s,A_t=a,S_{t+1}=s',R_{t+1}=r]
=\mathbb E_\pi[G_{t+1}\mid S_{t+1}=s'].
\]

Par définition de la fonction de valeur,
\[
\mathbb E_\pi[G_{t+1}\mid S_{t+1}=s'] = v_\pi(s').
\]

Donc
\[
\mathbb E_\pi[R_{t+1}+\gamma G_{t+1}\mid S_t=s,A_t=a,S_{t+1}=s',R_{t+1}=r]
= r + \gamma v_\pi(s').
\]

En revenant à la somme précédente,
\[
\mathbb E_\pi[R_{t+1}+\gamma G_{t+1}\mid S_t=s,A_t=a]
= \sum_{s',r} p(s',r\mid s,a)\big(r+\gamma v_\pi(s')\big).
\]

\section*{6. Équation de Bellman}

On réinjecte cette expression dans la formule obtenue à l'étape 4.1 :
\[
v_\pi(s)
=\sum_a \pi(a\mid s)
\sum_{s',r} p(s',r\mid s,a)\big(r+\gamma v_\pi(s')\big).
\]

Nous avons donc dérivé l'équation de Bellman pour la fonction de valeur d'état :
\[
\boxed{
v_\pi(s)
=\sum_a \pi(a\mid s)
\sum_{s',r} p(s',r\mid s,a)\big(r+\gamma v_\pi(s')\big)
}
\]

\section*{7. Version compacte avec une espérance imbriquée}

La même équation peut s'écrire de façon plus compacte :
\[
v_\pi(s)
=\mathbb E_\pi\big[R_{t+1}+\gamma v_\pi(S_{t+1})\mid S_t=s\big].
\]

Cette forme est exactement équivalente à la forme développée avec les sommes. Les deux écritures se déduisent l'une de l'autre par la définition de l'espérance d'une variable aléatoire discrète.

\section*{8. Lecture conceptuelle de chaque terme}

L'équation de Bellman dit que la valeur d'un état $s$ est :
\begin{enumerate}[leftmargin=2em]
    \item la moyenne sur les actions choisies par la politique $\pi(\cdot\mid s)$ ;
    \item la récompense immédiate moyenne ;
    \item plus la valeur actualisée du prochain état $\gamma v_\pi(s')$ ;
    \item le tout moyenné sur les transitions possibles du MDP.
\end{enumerate}

\section*{9. Remarque de rigueur}

La dérivation repose sur quatre ingrédients fondamentaux :
\begin{itemize}[leftmargin=2em]
    \item la décomposition algébrique du retour $G_t = R_{t+1}+\gamma G_{t+1}$ ;
    \item la linéarité de l'espérance conditionnelle ;
    \item la loi de l'espérance totale ;
    \item la propriété de Markov, qui permet de remplacer l'espérance conditionnelle du retour futur par $v_\pi(s')$ dès que l'état suivant est connu.
\end{itemize}

Sans la propriété de Markov, on ne pourrait pas, en général, résumer le futur uniquement par l'état $S_{t+1}$.

\end{document}
